\section{延伸：纽结同调的范畴化}\label{sec:categorification}

\subsection{这一章要解决什么}
本文把 Jones 多项式实现为三维理论的配分函数。一个自然的问题是：\textbf{Jones 多项式本身
是不是某个更丰富结构的影子？} 2000 年 Khovanov 给出了肯定的答案：Jones 多项式是一个
双分次同调群的 Euler 特征。本章说明这个"范畴化"（categorification）是什么、它与本文
是什么关系，以及为什么\textbf{不能}说本文预言了它。

\subsection{原文说了什么}
原文没有。全文 grep "homology"、"chain complex"、"categorif" 均零命中；
唯一擦边的是 p.372 关于"量子群"的一句，且带双重 hedging（\emph{should}、\emph{deferring a
fuller treatment}）\src{371--372}。本文的量子 Hilbert 空间取值在向量空间范畴，
没有 $\mathbb Z$-分次，也没有微分。

\subsection{Khovanov 构造：数 $\to$ 复形}
取链环 $L$ 的投影图，有 $n$ 个交叉。每个交叉有两种光滑化（0-光滑化与 1-光滑化），
于是有 $2^n$ 个顶点；每个顶点是一个平面图，配一个向量空间 $V^{\otimes k}$（$k$ 为圆分支数，
$V$ 是二维分次空间，$\dim V=2$，分次为 $q+q^{-1}$）。边对应鞍点（saddle）映射，
微分 $d$ 沿"改变一个光滑化"的边求和，满足 $d^2=0$。同调群 $\mathrm{Kh}(L)$ 双分次 $(i,j)$。
\begin{keypoint}[Euler 特征退回 Jones]
按分次取 Euler 特征：
\begin{equation}
\sum_{i,j}(-1)^i q^j\dim \mathrm{Kh}^{i,j}(L)=J(L;q). \tag{11.1}
\end{equation}
所以 Khovanov 同调是 Jones 多项式的"提升"：多项式是数，同调是复形。
注意归一化约定（$q$ 的幂次平移）因作者而异，比较前先对齐。
\end{keypoint}
以三叶结为例：$J=-q^{-4}+q^{-3}+q^{-1}$；Khovanov 同调的总维数是 $5$（$h$-维数），
分次维数表把每一项对应到 $J$ 的单项式。练习 1 就是查表核对 (11.1)。

\subsection{它比 Jones 多了什么}
\begin{enumerate}[leftmargin=*]
  \item \textbf{函子性}：给链环的协边（cobordism）配一个同调群之间的映射，
        得到函子 $\mathrm{Cob}\to\mathrm{Vect}$。Jones 多项式没有这种函子性。
        本文只到 p.397 的 "It is obvious at this point…"，连函子性都没写全。
  \item \textbf{解释整数性}：Jones 多项式的系数是整数，这在组合定义里是奇迹；
        在同调里是 Euler 特征的自然结果（Witten 2011 \S1.1.2 标题就叫
        "What Chern--Simons Theory Doesn't Explain"）。
  \item \textbf{更强的分辨力}：Khovanov 同调能区分一些 Jones 多项式相同的纽结
        （例如某些突变结）。
\end{enumerate}

\subsection{物理来源：为什么要"高一维"}
\begin{itemize}
  \item Witten 2011 把 Khovanov 同调实现为：四维 $\cN=4$ 超对称 Yang--Mills 在
        $W\times\mathbb R_+$ 上、以三维 CS 理论为边界条件；纽结对应边界上的 Wilson 算子。
        点题句：范畴化 "derived from a theory in one dimension higher"（W11 p.107）。
  \item Kapustin--Witten 2007 提供了四维扭曲理论的框架（几何 Langlands），
        但它本身不是 Khovanov 的推论；W11 引其 \S2、\S3、\S6.2。
  \item 换言之：本文的 CS 理论是三维的；范畴化需要四维的"母理论"。
        本文没有、也不可能预言这一点。
\end{itemize}

\subsection{要把本文理论范畴化，缺什么}
\begin{center}
\begin{tabular}{@{}p{3.2cm}p{8.8cm}@{}}
\toprule
缺口 & 说明 \\
\midrule
一维提升 & 需要一个四维理论，其边界是本文的 CS 理论（W11 用 4d $\cN=4$ SYM） \\
$Q^2=0$ 与双分次 & 需要超对称/BRST 结构给出微分与两个分次（$(i,j)$） \\
标架与半形式 & 反常必须典范化，否则只能得到射影对象（第 \ref{sec:surgery} 讲的相位） \\
表示范畴升 2-范畴 & $R$ 从向量空间升为链复形，融合从 $\otimes$ 升为 $\otimes^L$ \\
$q$ 的错位 & 本文 $q=e^{2\pi\ii/(k+h^\vee)}$ 是单位根；Khovanov 的 $q$ 是形式变量 \\
\bottomrule
\end{tabular}
\end{center}

\subsection{与本文的三个接口}
\begin{enumerate}[leftmargin=*]
  \item \textbf{配分函数值}：本文给 $Z(M;\text{links})\in\mathbb C$（数），范畴化给复形；
        取 Euler 特征回到数。这是"数 $\to$ 复形"的退化关系。
  \item \textbf{Wilson 线}：本文的 Wilson 线是三维算符；范畴化需要它成为四维理论的边界算子
        （W11 \S1.4.1）。本文的 Wilson 线是"终点"，范畴化是"起点"。
  \item \textbf{$q$ 与 level}：本文的 $q$ 由 $k$ 固定（单位根），Khovanov 的 $q$ 自由；
        两者只在 $q$ 取单位根时有对应，这一对应本身仍是研究课题。
\end{enumerate}

\subsection{练习}
\begin{enumerate}[leftmargin=*]
  \item 对三叶结（$J=-q^{-4}+q^{-3}+q^{-1}$），查 Khovanov 同调的分次维数表，
        验证 (11.1) 的 Euler 特征与多项式一致（注意归一化）。
  \item 说明为什么"双分次"比"单分次"更强：给出一个单分次无法区分、双分次可以区分的例子
        （提示：两分量链环）。
  \item 用本文的语言解释：为什么 $\cH_\Sigma$ 取值在 $\mathrm{Vect}$ 而不是链复形范畴？
  \item 把 p.372 那句量子群的话与 W11 的"高一维"对照，说明本文缺的是哪一维。
\end{enumerate}

\begin{boundary}
\begin{itemize}
  \item 已证（后世）：Khovanov 同调的 Euler 特征 = Jones 多项式；函子性（协边）。
  \item 未证/开放：Khovanov 同调的辛几何解释；与本文 CS 理论的直接函子关系。
  \item 本讲义立场：本章是延伸，不改变前十一讲的结论；任何"Witten 1989 预言了范畴化"的说法都是误读。
\end{itemize}
\end{boundary}
